\centerline{\bf \Huge Двойные отношения на окружности}

\begin{flushright}
        \scriptsize{мда треш}
\end{flushright}

\textbf{Определение.} Обобщенной окружностью на проективной плоскости называется прямая (само собой вместе с соответствующей ей бесконечно удаленной точкой) или же окружность.

\problem{1}
Точки $A, B, C, D$ лежат на обобщённой окружности. При инверсии, центр которой отличен от этих точек, они перешли в точки $A_1, B_1, C_1$, $D_1$ соответственно. Тогда $(A, B ; C, D)=\left(A_1, B_1 ; C_1, D_1\right)$. 

\textbf{Замечание.} Вообще говоря, это довольно известный факт из теории функций комплексной переменной, который звучит так: всякое \hyperlink{https://ru.wikipedia.org/wiki/Дробно-линейное_преобразование_комплексной_плоскости}{\textit{дробно-линейное}} преобразование \hyperlink{https://old.mccme.ru/ium//postscript/f15/geometry-02-linfrac.pdf}{сохраняет двойное отношение}.

\problem{2}
Пусть $A B C D$ - вписанный четырехугольник, в котором биссектрисы углов $A$ и $C$ пересекаются на диагонали $B D$. Докажите, что биссектрисы углов $B$ и $D$ пересекаются на диагонали $A C$.

\problem{3}
В окружности $S$ проведены две параллельные хорды $A B$ и $C D$. Прямая, проведенная через $C$ и середину $A B$, вторично пересекает $S$ в точке $E$. Точка $K$ - середина отрезка $D E$. Докажите, что $\angle A K E=\angle B K E$.

\problem{4}
Из точки $P$ к окружности $\omega$ проведены отрезки касательных $P A, P B$, точка $C$ диаметрально противоположна точке $B$. Докажите, что прямая $C P$ делит пополам перпендикуляр, опущенный из точки $A$ на прямую $B C$.

\problem{5}
В угол $B A C$ вписана окружность $\omega$, касающаяся сторон угла в точках $B, C$. Хорда $C D$ окружности $\omega$ параллельна прямой $A B$. Прямая $A D$ второй раз пересекает окружность $\omega$ в точке $E$. Докажите, что прямая $C E$ делит отрезок $A B$ пополам.

\problem{6}
$B_1$ и $B_2$ - основания внутренней и внешней биссектрис треугольника $A B C$. Касательные, отличные от прямой $A C$, проведенные из $B_1$ и $B_2$ к вписанной окружности, касаются её в точках $K_1$ и $K_2$. Докажите, что $B, K_1$ и $K_2$ лежат на одной прямой.

\problem{7}
Две неравные окружности $\omega_1$ и $\omega_2$ касаются внутренним образом окружности $\omega$ в точках $A$ и $B$. Пусть $C$ и $D$ точки пересечения окружностей $\omega_1$ и $\omega_2$. Прямая $C D$ пересекает $\omega$ в точках $E$ и $F$. Докажите, что касательные к $\omega$, проведенные в точках $E$ и $F$, пересекаются на прямой $A B$.

\newpage

\hspace{3mm}

\problem{8}
Четырехугольник $A B C D$ вписан в окружность $\omega$ с центром $O$. Биссектриса угла $A B D$ пересекает отрезок $A D$ в точке $K$ и окружность $\omega$ второй раз в точке $M$. Биссектриса угла $C B D$ пересекает отрезок $C D$ в точке $L$ и окружность $\omega$ второй раз в точке $N$. Известно, что прямые $K L$ и $M N$ параллельны. Докажите, что описанная окружность треугольника $M O N$ проходит через середину отрезка $B D$.


\problem{9}
Биссектриса внутреннего угла при вершине $A$ треугольника $A B C$ пересекает $(A B C)$ в точке $D$, а биссектриса внешнего угла пересекает $B C$ в точке $E$. Симедиана, поведённая из вершины $A$, пересекает $(A B C)$ в точке $L$. Докажите, что $D, E, L$ лежат на одной прямой.

\problem{10}
Окружности $\omega_1, \omega_2, \omega_3, \omega_4$ касаются друг друга внешним образом по циклу. Окружность $\omega$ касается их всех внешним образом в точках $X_1, X_2, X_3, X_4$ соответственно. Докажите, что $X_1 X_2 X_3 X_4$ образуют гармонический четырехугольник.

\problem{11}
Дан треугольник $A B C$ и точка $M$. Прямая, проходящая через $M$, пересекает $A B, B C, C A$ в $C_1, A_1, B_1$ соответственно. Прямые $A M, B M, C M$ пересекают окружность $(A B C)$ в точках $A_2, B_2, C_2$ соответственно. Докажите, что $A_1 A_2$, $B_1 B_2, C_1 C_2$ пересекаются в одной точке, лежащей на ( $A B C$ ).

\problem{12}
\textbf{Теорема о двойной бабочке.}
На хорде $M N$ окружности $\omega$ отмечены точки $M_1$ и $N_1$ такие, что $M M_1=N N_1$. Хорда $A C$ проходит через точку $M_1$, а хорда $B D$ - через точку $N_1$. Отрезки $A D$ и $B C$ пересекают $M N$ в точках $X$ и $Y$ соответственно. Докажите, что $X M=Y N$.

\problem{13}
Из точки $P$ вне окружности $\omega$ проведены касательные $P X$ и $P Y$ и две секущие, пересекающие окружность в точках $A$ и $B, C$ и $D$. Докажите, что точки пересечения прямых $A C$ и $B D, A D$ и $B C$ лежат на прямой $X Y$.

\problem{14}
Точки $D$ и $E$ - проекции вершин $B$ и $C$ на биссектрису угла $A$ треугольника $A B C$. Окружность, построенная на $D E$ как на диаметре, пересекает сторону $B C$ в точках $X$ и $Y$. Докажите, что $\angle B A X=\angle C A Y$.

\problem{15}
Окружность проходит через вершины $B$ и $C$ треугольника $A B C$ и повторно пересекает отрезки $A B$ и $A C$ в точках $C_1$ и $B_1$ соответственно. Отрезки $B B_1$ и $C C_1$ пересекаются в точке $P$, прямая $A P$ пересекает отрезок $B C$ в точке $A_1$. Докажите, что окружность $(A_1 B_1 C_1)$ проходит через середину отрезка $B C$.